\INDEX{證明：}首先，目前仍未知直線$l$與$\triangle{BEF}$的外接圓是否存在交點，如果要用「證明交點為切點」的方法來證明，
必須先證明兩者之間存在交點，這個方法只會增加證明的難度。

既然如此，那麼更加不可能使用「切割線定理」來證明命題，因為切割線定理的前提是已知兩者相交。

經過篩選，目前可以使用的證明方法就有「直線$l$與$\triangle{BEF}$外接圓圓心的距離$d_1$等於此外接圓半徑$R_1$」。
\begin{center} % 子圖 1
  \begin{tikzpicture}[scale=1.2]
    % --- 1. 定義圓心 M, N 與半徑 ---
    \tkzDefPoint(-2,0){M}
    \tkzDefPoint(3,0){N}
    \def\rO{2}
    \def\rG{4}
    
    % --- 2. 定義直線 MN 上的 C, D (確保 C, M, N, D 順序) ---
    \tkzDefPoint(-2-\rO,0){C}
    \tkzDefPoint(3+\rG,0){D}
    
    % --- 3. 繪製兩圓 Omega, Gamma ---
    %\tkzDrawCircle[thick,black](M,C)
    %\tkzDrawCircle[thick,black](N,D)
    
    % --- 4. 求兩圓交點 A, B ---
    \tkzInterCC(M,C)(N,D) \tkzGetPoints{A}{B}
    
    % --- 5. 定義 P 為三角形 ACD 的外心 ---
    \tkzCircumCenter(A,C,D) \tkzGetPoint{P}
    %\tkzDrawCircle[dashed, cyan](P,A) % 繪製 ACD 的外接圓供參考
    %\tkzDrawArc[color=red, dashed](P,D)(C)
    % --- 6. 構造割線 AP 並求交點 E, F ---
    \tkzInterLC(A,P)(M,C) \tkzGetPoints{E}{A_tmp}
    \tkzInterLC(A,P)(N,D) \tkzGetPoints{A_tmp2}{F}
    
    % --- 7. 計算三角形 PMN 的垂心 H ---
    \tkzOrthoCenter(P,M,N) \tkzGetPoint{H}
    
    % --- 8. 繪製三角形 BEF 的外接圓 ---
    \tkzCircumCenter(B,E,F) \tkzGetPoint{O_1}
    \tkzDrawCircle[dashed, black](O_1,B)
    
    % --- 9. 繪製通過 H 且平行於 AP 的切線 l ---
    \tkzDefLine[parallel=through H](A,P) \tkzGetPoint{l_dir}
    \tkzDrawLine[blue, thick, add=0.1 and -0.7](H,l_dir)
    
    % --- 10. 繪製線段與標籤 ---
    %\tkzDrawLines[add=0.1 and 0.1](C,D)
    %\tkzDrawSegments(A,P P,M P,N A,B B,E B,F)
    %\tkzDrawPoints[fill=black](M,N,A,B,C,D,E,F,P,H)
    \tkzDrawPoints[fill=black](B,E,F,H,O_1)
    %\tkzLabelPoints[above](M,N)
    %\tkzLabelPoints[above left](C)
    %\tkzLabelPoints[above right](D,H)
    \tkzLabelPoints[above right](H)
    \tkzLabelPoints[left](B)
    %\tkzLabelPoints[above](A)
    %\tkzLabelPoints[below](P)
    \tkzLabelPoints[above](E)
    \tkzLabelPoints[below](F)
    \tkzLabelPoints[above](O_1)
    
    %\node at (-2-\rO*0.8,\rO*0.8) {$\Omega$};
    %\node at (3+\rG*0.8,\rG*0.8) {$\Gamma$};
    
    \coordinate (l_1) at ($(H)!0.2!(l_dir)$);
    \node[blue, right] at (l_1) {$l$};
    % 1. 找出 O1 在直線 (H, l_dir) 上的投影點（垂足），命名為 K
    \tkzDefPointBy[projection = onto H--l_dir](O_1) \tkzGetPoint{K}

    % 2. 繪製垂線段 O1K
    \tkzDrawSegment[color=red, dashed](O_1,K)

    % 3. 繪製直角符號
    \tkzMarkRightAngle[size=0.2](O_1,K,l_dir)

    \coordinate (d_1) at ($(O_1)!0.5!(K)$);
    \node[red, below] at (d_1) {$d_1$};
    \tkzDrawSegments(B,E B,F E,F)

     % 1. 找出 H 在直線 (E, F) 上的投影點（垂足），命名為 M
    \tkzDefPointBy[projection = onto E--F](H) \tkzGetPoint{M}

    % 2. 繪製垂線段 O1K
    \tkzDrawSegment[color=purple, dashed](H,M)

    % 3. 繪製直角符號
    \tkzMarkRightAngle[size=0.2](H,M,F)
    \coordinate (d_2) at ($(H)!0.5!(M)$);
    \node[purple, below right] at (d_2) {$d_2$};
  \end{tikzpicture}
\end{center}
現在把幾何圖的子圖放大，並記一對平行線$\left(EF,l\right)$之間的距離為$d_2$，則原命題等價於下方的命題一：
\begin{equation}
  d_1=R_1\Leftrightarrow d_2=R_1+\left\{O_1\text{到}EF\text{的距離}\right\}=R_1-R_1\cos{\angle{EBF}}
\end{equation}
同時，$d_2$ 的長度等於 $H$ 到 $EF$ 的距離。另一方面，$H$為$\triangle{BEF}$的垂心，且$P$在直線$EF$上，所以連結$HP$是必然的。
\begin{center} % 子圖 2
  \begin{tikzpicture}[rotate=75, scale=1.2]
    % --- 1. 定義圓心 M, N 與半徑 ---
    \tkzDefPoint(-2,0){M}
    \tkzDefPoint(3,0){N}
    \def\rO{2}
    \def\rG{4}
    
    % --- 2. 定義直線 MN 上的 C, D (確保 C, M, N, D 順序) ---
    \tkzDefPoint(-2-\rO,0){C}
    \tkzDefPoint(3+\rG,0){D}
    
    % --- 3. 繪製兩圓 Omega, Gamma ---
    %\tkzDrawCircle[thick,black](M,C)
    %\tkzDrawCircle[thick,black](N,D)
    
    % --- 4. 求兩圓交點 A, B ---
    \tkzInterCC(M,C)(N,D) \tkzGetPoints{A}{B}
    
    % --- 5. 定義 P 為三角形 ACD 的外心 ---
    \tkzCircumCenter(A,C,D) \tkzGetPoint{P}
    %\tkzDrawCircle[dashed, cyan](P,A) % 繪製 ACD 的外接圓供參考
    %\tkzDrawArc[color=red, dashed](P,D)(C)
    % --- 6. 構造割線 AP 並求交點 E, F ---
    \tkzInterLC(A,P)(M,C) \tkzGetPoints{E}{A_tmp}
    \tkzInterLC(A,P)(N,D) \tkzGetPoints{A_tmp2}{F}
    
    % --- 7. 計算三角形 PMN 的垂心 H ---
    \tkzOrthoCenter(P,M,N) \tkzGetPoint{H}
    
    % --- 8. 繪製三角形 BEF 的外接圓 ---
    \tkzCircumCenter(B,E,F) \tkzGetPoint{O_1}
    \tkzDrawCircle[dashed, black](O_1,B)
    
    % --- 9. 繪製通過 H 且平行於 AP 的切線 l ---
    \tkzDefLine[parallel=through H](A,P) \tkzGetPoint{l_dir}
    \tkzDrawLine[blue, thick, add=0.1 and -0.7](H,l_dir)
    
    % --- 10. 繪製線段與標籤 ---
    %\tkzDrawLines[add=0.1 and 0.1](C,D)
    %\tkzDrawSegments(A,P P,M P,N A,B B,E B,F)
    %\tkzDrawPoints[fill=black](M,N,A,B,C,D,E,F,P,H)
    \tkzDrawPoints[fill=black](B,E,F,H,P)
    %\tkzLabelPoints[above](M,N)
    %\tkzLabelPoints[above left](C)
    %\tkzLabelPoints[above right](D,H)
    \tkzLabelPoints[above](H)
    \tkzLabelPoints[below](B)
    %\tkzLabelPoints[above](A)
    \tkzLabelPoints[below](P)
    \tkzLabelPoints[left](E)
    \tkzLabelPoints[below right](F)
    %\tkzLabelPoints[above](O_1)
    
    %\node at (-2-\rO*0.8,\rO*0.8) {$\Omega$};
    %\node at (3+\rG*0.8,\rG*0.8) {$\Gamma$};
    
    \coordinate (l_1) at ($(H)!0.2!(l_dir)$);
    \node[blue, above] at (l_1) {$l$};
    % 1. 找出 O1 在直線 (H, l_dir) 上的投影點（垂足），命名為 K
    %\tkzDefPointBy[projection = onto H--l_dir](O_1) \tkzGetPoint{K}

    % 2. 繪製垂線段 O1K
    %\tkzDrawSegment[color=red, dashed](O_1,K)

    % 3. 繪製直角符號
    %\tkzMarkRightAngle[size=0.2](O_1,K,l_dir)

    %\coordinate (d_1) at ($(O_1)!0.5!(K)$);
    %\node[red, below] at (d_1) {$d_1$};
    \tkzDrawSegments(B,E B,F E,P)
    \tkzDrawSegments[dashed](H,P)

     % 1. 找出 H 在直線 (E, F) 上的投影點（垂足），命名為 M
    \tkzDefPointBy[projection = onto E--F](H) \tkzGetPoint{M}

    % 2. 繪製垂線段 O1K
    \tkzDrawSegment[color=purple, dashed](H,M)

    % 3. 繪製直角符號
    \tkzMarkRightAngle[size=0.2](H,M,F)
    \coordinate (d_2) at ($(H)!0.5!(M)$);
    \node[purple, below right] at (d_2) {$d_2$};

\end{tikzpicture}
\end{center}
所以命題一等價於下方的命題二：
\begin{equation}
  HP\sin{\angle{HPE}}=d_2=R_1\left(1-\cos{\angle{EBF}}\right)
\end{equation}
現在只剩下$R_1$仍未計算出來，所以下一步就是計算$R_1$的長度。對$\triangle{BEF}$使用正弦定理可得：
\begin{equation}
  2R_1=\frac{BE}{\sin{\angle{BFE}}}\left(=\frac{BF}{\sin{\angle{BEF}}}\right)
\end{equation}
所以命題二等價於下方的命題三：
\begin{equation}
  HP\sin{\angle{HPE}}=\frac{BE}{2\sin{\angle{BFE}}}\left(1-\cos{\angle{EBF}}\right)
\end{equation}

下一步，我們需要把命題三中的各個量轉化成與$E,F$無關的量：
\begin{itemize}
  \item 因為$A,B$是圓$\Omega$與$\Gamma$的交點，所以$AB$是兩圓的公共弦，因此$AB\perp MN$。
  同時，$H$是$\triangle{PMN}$的垂心，所以$PH\perp MN$。由此可得 $AB\parallel PH$，所以$$\angle{HPE}=\angle{BAP}$$
  \item 記圓$\Omega$的半徑為$R_M$，對圓$\Omega$使用正弦定理：$$BE=2R_M\sin{\angle{BAE}}=2R_M\sin{\angle{BAP}}$$
  \item 在圓$\Gamma$中使用Angle Chasing：$\angle{BFE}=\angle{BFA}=\angle{BDA}$。記圓$\Gamma$的半徑為$R_N$，對圓$\Gamma$使用正弦定理：
  $$\sin{\angle{BFE}}=\sin{\angle{BDA}}=\frac{BA}{2R_N}$$
  \item $$\angle{EBF}=180^{\circ}-\angle{BEF}-\angle{BFE}=180^{\circ}-\angle{BEF}-\angle{BDA}$$
  在圓$\Gamma$中使用Angle Chasing：$$\angle{BEF}=\angle{BCA}$$
  所以有$$\angle{EBF}=180^{\circ}-\angle{BCA}-\angle{BDA}$$
  \iffalse
  利用圓周角定理可知
  $$\begin{cases}\angle{BCA}=\frac{1}{2}\angle{BMA}=\angle{AMN}\\ \angle{BDA}=\frac{1}{2}\angle{BNA}=\angle{ANM}\end{cases}$$
  因此
  $$\angle{EBF}=180^{\circ}-\angle{AMN}-\angle{ANM}=\angle{MAN}$$
  \fi
\end{itemize}
將上述結果代入命題三，可得命題四：
\begin{align}
  HP\sin{\angle{BAP}}&=\frac{2R_M\sin{\angle{BAP}}}{2\cdot\frac{BA}{2R_N}}
  \left[1-\cos{\left(180^{\circ}-\angle{BCA}-\angle{BDA}\right)}\right]\notag\\
  &=\frac{2R_MR_N\sin{\angle{BAP}}}{AB}\left[1+\cos{\left(\angle{BCA}+\angle{BDA}\right)}\right]
\end{align}
命題四兩邊都存在$\sin{\angle{BAP}}$，所以可以約去，得到命題五：
\begin{equation}
  HP=\frac{2R_MR_N}{AB}\left[1+\cos{\left(\angle{BCA}+\angle{BDA}\right)}\right]
\end{equation}
下一步，我們需要想辦法把$HP$的長度表示成與$H$無關的量，所以我把另一個幾何圖的子圖放大
\begin{center} % 子圖 3
  \begin{tikzpicture}[rotate=-90, scale=0.6]
    \tkzDefPoint(-2,0){M}
    \tkzDefPoint(3,0){N}
    \def\rO{2}
    \def\rG{4}
    \tkzDefPoint(-2-\rO,0){C}
    \tkzDefPoint(3+\rG,0){D}
    \tkzInterCC(M,C)(N,D) \tkzGetPoints{A}{B}
    \tkzCircumCenter(A,C,D) \tkzGetPoint{P}
    \tkzInterLC(A,P)(M,C) \tkzGetPoints{E}{A_tmp}
    \tkzInterLC(A,P)(N,D) \tkzGetPoints{A_tmp2}{F}
    \tkzOrthoCenter(P,M,N) \tkzGetPoint{H}
    \tkzDrawSegments(M,N P,N P,M)
    \tkzDrawPoints[fill=black](M,N,P,H)    
    \tkzLabelPoints[right](M,N)
    \tkzLabelPoints[above left](H)
    \tkzLabelPoints[below](P)
  \end{tikzpicture}
\end{center}
如果大家還記得「垂心-重心-外心」關係的話，很快就能得到：
$$
\overrightarrow{O_2H}=\overrightarrow{O_2P}+\overrightarrow{O_2M}+\overrightarrow{O_2N}
\Leftrightarrow\overrightarrow{PH}=\overrightarrow{O_2M}+\overrightarrow{O_2N}
$$
其中$O_2$為$\triangle{PMN}$的外心，為了讀者能更清晰地去計算，下圖為添加了$O_2$的幾何圖子圖：
\begin{center} % 子圖 4
  \begin{tikzpicture}[rotate=-90, scale=0.6]
    \tkzDefPoint(-2,0){M}
    \tkzDefPoint(3,0){N}
    \def\rO{2}
    \def\rG{4}
    \tkzDefPoint(-2-\rO,0){C}
    \tkzDefPoint(3+\rG,0){D}
    \tkzInterCC(M,C)(N,D) \tkzGetPoints{A}{B}
    \tkzCircumCenter(A,C,D) \tkzGetPoint{P}
    \tkzInterLC(A,P)(M,C) \tkzGetPoints{E}{A_tmp}
    \tkzInterLC(A,P)(N,D) \tkzGetPoints{A_tmp2}{F}
    \tkzOrthoCenter(P,M,N) \tkzGetPoint{H}
    \tkzCircumCenter(P,M,N) \tkzGetPoint{O_2}
    \tkzDrawCircle[dashed, black](O_2,M)
    \tkzDrawSegments(M,N P,N P,M)
    \tkzDrawSegments[dashed, black](O_2,M O_2,N)
    \tkzDrawPoints[fill=black](M,N,P,H,O_2)   
    \tkzLabelPoints[right](M,N)
    \tkzLabelPoints[above left](H)
    \tkzLabelPoints[below](P)
    \tkzLabelPoints[below](O_2)
  \end{tikzpicture}
\end{center}
記此外接圓半徑為$R_2$，由幾何關係可以算得
\begin{align*}
  HP&=\left|\overrightarrow{O_2M}+\overrightarrow{O_2N}\right|=2R_2\cos{\frac{\angle{MO_2N}}{2}}\\
  &=2R_2\cos{\angle{MPN}}=\frac{MN}{\sin{}\angle{MPN}}\cos{\angle{MPN}}\\
  &=MN\cot{\angle{MPN}}
\end{align*}
代入命題五可得
\begin{equation}
  MN\cot{\angle{MPN}}=\frac{2R_MR_N}{AB}\left[1+\cos{\left(\angle{BCA}+\angle{BDA}\right)}\right]
\end{equation}
現在整理當前的幾何圖
\begin{center} % 子圖 5
  \begin{tikzpicture}[scale=0.8]
    \tkzDefPoint(-2,0){M}
    \tkzDefPoint(3,0){N}
    \def\rO{2}
    \def\rG{4} 
    \tkzDefPoint(-2-\rO,0){C}
    \tkzDefPoint(3+\rG,0){D}
    \tkzDrawCircle[thick,black](M,C)
    \tkzDrawCircle[thick,black](N,D)
    \tkzInterCC(M,C)(N,D) \tkzGetPoints{A}{B}
    \tkzCircumCenter(A,C,D) \tkzGetPoint{P}
    \tkzDrawArc[color=red, dashed](P,D)(C)
    \tkzDrawLines[add=0.1 and 0.1](C,D)
    \tkzDrawSegments(A,P P,M P,N A,B)
    \tkzDrawPoints[fill=black](M,N,A,B,C,D,P)
    \tkzLabelPoints[above](M,N)
    \tkzLabelPoints[above left](C)
    \tkzLabelPoints[above right](D)
    \tkzLabelPoints[above left](B)
    \tkzLabelPoints[above](A)
    \tkzLabelPoints[below](P)
    \node at (-2-\rO*0.8,\rO*0.8) {$\Omega$};
    \node at (3+\rG*0.8,\rG*0.8) {$\Gamma$};
  \end{tikzpicture}
\end{center}
下一步，我們需要把$\angle{MPN}$表示成與$P$無關的量，即把$P$點消去：

因為$P$為$\triangle{ACD}$的外心，而$M$為$\triangle{ACB}$的外心，所以直線$MP$為$AC$的中垂線。
同理，直線$NP$為$AD$的中垂線。

為了能夠更清楚地解釋，筆者再提取了幾何圖的子圖：
\begin{center} % 子圖 6
  \begin{tikzpicture}[scale=0.6]
    \tkzDefPoint(-2,0){M}
    \tkzDefPoint(3,0){N}
    \def\rO{2}
    \def\rG{4}
    \tkzDefPoint(-2-\rO,0){C}
    \tkzDefPoint(3+\rG,0){D}
    \tkzDrawCircle[thick,black](M,C)
    \tkzDrawCircle[thick,black](N,D)
    \tkzInterCC(M,C)(N,D) \tkzGetPoints{A}{B}
    \tkzCircumCenter(A,C,D) \tkzGetPoint{P}
    \tkzDrawArc[color=red, dashed](P,D)(C)
    \tkzInterLL(M,P)(A,C) \tkzGetPoint{X}
    \tkzInterLL(N,P)(A,D) \tkzGetPoint{Y}
    \tkzDrawSegments(A,P P,M P,N)
    \tkzDrawSegments[dashed](M,X A,C N,Y A,D)
    \tkzDrawPoints[fill=black](M,N,A,C,D,P,X,Y)
    \tkzLabelPoints[below left](M)
    \tkzLabelPoints[below right](N)
    \tkzLabelPoints[above left](C)
    \tkzLabelPoints[above right](D)
    \tkzLabelPoints[above](A)
    \tkzLabelPoints[below](P)
    \tkzLabelPoints[above](X,Y)
    \tkzMarkRightAngle[size=0.2](M,X,A)
    \tkzMarkRightAngle[size=0.2](N,Y,A)
  \end{tikzpicture}
\end{center}
記$PM$交$AC$於$X$，$PN$交$AD$於$Y$，由$PM\perp AC$和$PN\perp AD$可知$A,X,P,Y$四點共圖，所以
$$\angle{MPN}=\angle{XPY}=180^{\circ}-\angle{XAY}=180^{\circ}-\angle{CAD}$$
代入至命題五可得：
\begin{equation}
  \frac{2R_MR_N}{AB}\left[1+\cos{\left(\angle{BCA}+\angle{BDA}\right)}\right]
  =MN\cot{\left(180^{\circ}-\angle{CAD}\right)}=-MN\cot{\angle{CAD}}
\end{equation}
現在整理當前的幾何圖
\begin{center} % 子圖 7
  \begin{tikzpicture}[scale=0.8]
    \tkzDefPoint(-2,0){M}
    \tkzDefPoint(3,0){N}
    \def\rO{2}
    \def\rG{4}
    \tkzDefPoint(-2-\rO,0){C}
    \tkzDefPoint(3+\rG,0){D}
    \tkzDrawCircle[thick,black](M,C)
    \tkzDrawCircle[thick,black](N,D)
    \tkzInterCC(M,C)(N,D) \tkzGetPoints{A}{B}
    \tkzDrawLines[add=0.1 and 0.1](C,D)
    \tkzDrawSegments(A,B C,A A,D)
    \tkzDrawPoints[fill=black](M,N,A,B,C,D)
    \tkzLabelPoints[above](M,N)
    \tkzLabelPoints[above left](C)
    \tkzLabelPoints[above right](D)
    \tkzLabelPoints[above left](B)
    \tkzLabelPoints[above](A)
    \node at (-2-\rO*0.8,\rO*0.8) {$\Omega$};
    \node at (3+\rG*0.8,\rG*0.8) {$\Gamma$};
  \end{tikzpicture}
\end{center}
下一步，我們要消去$C,D$兩點：
利用圓周角定理可知
$$
\begin{cases}
  \angle{BCA}=\frac{1}{2}\angle{BMA}=\angle{AMN}\\ 
  \angle{BDA}=\frac{1}{2}\angle{BNA}=\angle{ANM}
\end{cases}
$$
所以有
$$
\angle{BCA}+\angle{BDA}=\angle{AMN}+\angle{ANM}=180^{\circ}-\angle{MAN}
$$
另一方面
\begin{align*}
  \angle{CAD}&=\angle{MAN}+\angle{CAM}+\angle{DAN}\\
  &=\angle{MAN}+\angle{MAC}+\angle{NDA}\\
  &=\angle{MAN}+\frac{1}{2}\left(\angle{BCA}+\angle{BDA}\right)\\
  &=\angle{MAN}+\frac{1}{2}\left(180^{\circ}-\angle{MAN}\right)\\
  &=90^{\circ}+\frac{1}{2}\angle{MAN}
\end{align*}
因此命題五等價於下方的命題六：
\begin{align}
  &\frac{2R_MR_N}{AB}\left[1+\cos{\left(180^{\circ}-\angle{MAN}\right)}\right]
  =-MN\cot{\left(90^{\circ}+\frac{1}{2}\angle{MAN}\right)}\notag\\
  \Leftrightarrow\quad&\frac{2R_MR_N}{AB}\left(1-\cos{\angle{MAN}}\right)
  =MN\tan{\frac{\angle{MAN}}{2}}\notag\\
  \Leftrightarrow\quad&\frac{2R_MR_N}{AB}\cdot2\sin^2{\frac{\angle{MAN}}{2}}=MN\cdot\tan{\frac{\angle{MAN}}{2}}\notag\\
  \Leftrightarrow\quad&\frac{2R_MR_N}{AB}\cdot2\sin{\frac{\angle{MAN}}{2}}\cos{\frac{\angle{MAN}}{2}}=MN\notag\\
  \Leftrightarrow\quad&\frac{2R_MR_N}{AB}\sin{\angle{MAN}}=MN\notag\\
  \Leftrightarrow\quad&2R_MR_N\sin{\angle{MAN}}=MN\cdot AB
\end{align}
現在整理當前的幾何圖
\begin{center} % 子圖 8
  \begin{tikzpicture}[scale=1]
    \tkzDefPoint(-2,0){M}
    \tkzDefPoint(3,0){N}
    \def\rO{2}
    \def\rG{4}
    \tkzDefPoint(-2-\rO,0){C}
    \tkzDefPoint(3+\rG,0){D}
    \tkzInterCC(M,C)(N,D) \tkzGetPoints{A}{B}
    \tkzDrawArc[thick,black](M,B)(A)
    \tkzDrawArc[thick,black](N,A)(B)
    \tkzDrawSegments(A,B A,M A,N B,M B,N M,N)
    \tkzDrawPoints[fill=black](M,N,A,B)
    \tkzLabelPoints[left](M)
    \tkzLabelPoints[right](N)
    \tkzLabelPoints[below](B)
    \tkzLabelPoints[above](A)
    \tkzInterLL(M,N)(A,B) \tkzGetPoint{K}
    \tkzMarkRightAngle[size=0.2, red](A,K,N)
  \end{tikzpicture}
\end{center}
因為 $MN\perp AB$，所以
\begin{equation}
  \left[AMBN\right]=\frac{1}{2}MN\cdot AB\label{eq:Area_AMBN_1}
\end{equation}
另一方面，易知 $\triangle{AMN}\cong\triangle{BMN}$，所以
\begin{align}
  \left[AMBN\right]&=\left[\triangle{AMN}\right]+\left[\triangle{BMN}\right]\notag\\
  &=2\left[\triangle{AMN}\right]\notag\\
  &=AM\cdot AN\sin{\angle{MAN}}\notag\\
  &=R_MR_N\sin{\angle{MAN}}\label{eq:Area_AMBN_2}
\end{align}
聯立式\eqref{eq:Area_AMBN_1}與式\eqref{eq:Area_AMBN_2}可得
\begin{align}
  &\frac{1}{2}MN\cdot AB=R_MR_N\sin{\angle{MAN}}\notag\\
  \Leftrightarrow\quad&2R_MR_N\sin{\angle{MAN}}=MN\cdot AB
\end{align}
由此可證得命題六成立，即原命題成立。